% Compatibilidad con instalaciones MiKTeX cuyo formato base es anterior a
% los paquetes Beamer/Hyperref instalados.
\providecommand{\IfFormatAtLeastTF}[3]{#3}
\providecommand{\IfFormatAtLeastT}[2]{}
\providecommand{\IfFormatAtLeastF}[2]{#2}
\providecommand{\IfDocumentMetadataTF}[2]{#2}
\providecommand{\IfDocumentMetadataT}[1]{}
\providecommand{\IfDocumentMetadataF}[1]{#1}
\makeatletter
\@namedef{ver@epstopdf-base.sty}{9999/12/31}
\makeatother

\documentclass[aspectratio=169,11pt]{beamer}

% =============================================================================
% MACROECONOMIA II | LICENCIATURA EN ECONOMIA | CIDE
% SESION 5: MODELO INTERTEMPORAL REAL CON INVERSION
% =============================================================================

\usepackage[utf8]{inputenc}
\usepackage[english]{babel}
\uselanguage{English}
\languagepath{English}
\usepackage{graphicx}
\usepackage{amsmath,amssymb}
\IfFileExists{booktabs.sty}{
  \usepackage{booktabs}
}{
  \providecommand{\toprule}{\hline}
  \providecommand{\midrule}{\hline}
  \providecommand{\bottomrule}{\hline}
}
\providecommand{\addlinespace}[1][0.5em]{}
\usepackage{xcolor}
\usepackage{ragged2e}
\usepackage{etoolbox}
\usepackage{array,tabularx}

\newtheorem{proposition}[theorem]{Proposición}
\newtheorem{remark}[theorem]{Observación}
\newtheorem{assumption}[theorem]{Supuesto}

\AtBeginDocument{\justifying}
\patchcmd{\itemize}{\raggedright}{\justifying}{}{}
\patchcmd{\enumerate}{\raggedright}{\justifying}{}{}
\patchcmd{\description}{\raggedright}{\justifying}{}{}
\addtobeamertemplate{block begin}{}{\justifying}
\addtobeamertemplate{block alerted begin}{}{\justifying}
\addtobeamertemplate{block example begin}{}{\justifying}
\makeatletter
\patchcmd{\beamer@columnenv}{\raggedright}{\justifying}{}{}
\patchcmd{\beamer@columncom}{\raggedright}{\justifying}{}{}
\makeatother

\definecolor{CIDENavy}{RGB}{3,25,54}
\definecolor{CIDEBlue}{RGB}{5,35,75}
\definecolor{CIDEInk}{RGB}{25,34,45}
\definecolor{CIDEGray}{RGB}{101,112,126}
\definecolor{CIDEGrayLight}{RGB}{243,244,246}

\setbeamercolor{normal text}{fg=CIDEInk,bg=white}
\setbeamercolor{structure}{fg=CIDEBlue}
\setbeamercolor{alerted text}{fg=CIDEBlue}
\setbeamercolor{title}{bg=CIDEBlue,fg=white}
\setbeamercolor{subtitle}{fg=CIDEInk}
\setbeamercolor{frametitle}{bg=CIDEBlue,fg=white}
\setbeamercolor{palette primary}{bg=CIDEBlue,fg=white}
\setbeamercolor{palette secondary}{bg=CIDEBlue,fg=white}
\setbeamercolor{palette tertiary}{bg=CIDENavy,fg=white}
\setbeamercolor{palette quaternary}{bg=CIDENavy,fg=white}
\setbeamercolor{section in head/foot}{bg=CIDENavy,fg=white}
\setbeamercolor{section in head/foot shaded}{bg=CIDENavy,fg=white!50}
\setbeamercolor{mini frame}{bg=white,fg=white}
\setbeamercolor{mini frame shaded}{bg=white!35,fg=white!35}
\setbeamercolor{upper separation line head}{bg=CIDENavy}
\setbeamercolor{lower separation line head}{bg=CIDEBlue}
\setbeamercolor{block title}{bg=CIDEBlue,fg=white}
\setbeamercolor{block body}{bg=CIDEGrayLight,fg=CIDEInk}
\setbeamercolor{block title alerted}{bg=CIDENavy,fg=white}
\setbeamercolor{block body alerted}{bg=CIDEGrayLight,fg=CIDEInk}
\setbeamercolor{caption name}{fg=CIDEBlue}

\setbeamerfont{author in head/foot}{size=\tiny}
\setbeamerfont{institute in head/foot}{size=\tiny}
\setbeamerfont{title in head/foot}{size=\tiny}
\setbeamersize{text margin left=0.70cm,text margin right=0.70cm}
\setlength{\leftmargini}{1.2em}
\setlength{\leftmarginii}{1.1em}
\setlength{\abovecaptionskip}{0.25em}
\setlength{\belowcaptionskip}{0pt}

\useoutertheme[subsection=false]{miniframes}
\setbeamertemplate{navigation symbols}{}
\setbeamertemplate{mini frame}[circle]
\setbeamertemplate{mini frame in current subsection}[circle]
\setbeamertemplate{section in head/foot}{\color{white}\insertsectionhead}
\setbeamertemplate{section in head/foot shaded}{\color{white!55}\insertsectionhead}
\setbeamertemplate{caption}[numbered]
\setbeamertemplate{section in toc}[circle]
\setbeamertemplate{subsection in toc}[circle]
\setbeamertemplate{itemize item}{\raisebox{0.15ex}{\tiny$\blacktriangleright$}}
\setbeamertemplate{itemize subitem}{\raisebox{0.15ex}{\tiny$\bullet$}}
\setbeamertemplate{enumerate item}[circle]

\setbeamertemplate{frametitle}{%
  \nointerlineskip
  \begin{beamercolorbox}[
    wd=\paperwidth,ht=3.05ex,dp=1.05ex,
    leftskip=0.42cm,rightskip=0.42cm
  ]{frametitle}
    \usebeamerfont{frametitle}\insertframetitle
  \end{beamercolorbox}%
}

\setbeamertemplate{footline}{%
  \leavevmode
  \vbox{%
    \offinterlineskip
    \hbox{%
      \begin{beamercolorbox}[wd=.72\paperwidth,ht=2.05ex,dp=.55ex,leftskip=.35cm]{palette secondary}
        \usebeamerfont{author in head/foot}\insertshortauthor
      \end{beamercolorbox}%
      \begin{beamercolorbox}[wd=.28\paperwidth,ht=2.05ex,dp=.55ex,rightskip=.35cm plus1fil]{palette secondary}
        \hfill\usebeamerfont{institute in head/foot}\insertshortinstitute
      \end{beamercolorbox}%
    }%
    \hbox{%
      \begin{beamercolorbox}[wd=.86\paperwidth,ht=1.85ex,dp=.50ex,leftskip=.35cm]{palette tertiary}
        \usebeamerfont{title in head/foot}\insertshorttitle
      \end{beamercolorbox}%
      \begin{beamercolorbox}[wd=.14\paperwidth,ht=1.85ex,dp=.50ex,center]{palette tertiary}
        \usebeamerfont{title in head/foot}\insertframenumber{} / \inserttotalframenumber
      \end{beamercolorbox}%
    }%
  }%
  \vskip0pt
}

\newcommand{\RutaLogoCIDE}{Figuras/General/logocide.png}
\setbeamertemplate{background}{%
  \vbox to \paperheight{%
    \vfill
    \hbox to \paperwidth{%
      \hfill
      \includegraphics[height=0.58cm,keepaspectratio]{\RutaLogoCIDE}%
      \hspace{0.18cm}%
    }%
    \vspace{0.68cm}%
  }%
}

\setbeamertemplate{title page}{%
  \vspace*{0.35cm}
  \centering
  \begin{beamercolorbox}[wd=\textwidth,sep=0.24cm,center,rounded=false]{title}
    \usebeamerfont{title}\inserttitle\par
  \end{beamercolorbox}
  \vspace{0.25cm}
  {\usebeamerfont{subtitle}\usebeamercolor[fg]{subtitle}\insertsubtitle\par}
  \vfill
  {\usebeamerfont{author}\insertauthor\par}
  \vspace{0.35cm}
  {\usebeamerfont{institute}\insertinstitute\par}
  \vfill
  {\usebeamerfont{date}\insertdate\par}
  \vspace{0.25cm}
}

\newcommand{\FigurasSesion}{Figuras/05_Modelo Intertemporal con Inversion}
\newcommand{\figura}[2][0.65\textheight]{%
  \includegraphics[width=\linewidth,height=#1,keepaspectratio]{\FigurasSesion/#2}%
}
\newcommand{\cidehighlight}[1]{\textcolor{CIDEBlue}{\bfseries #1}}
\newcommand{\fuentefigura}[1]{%
  \vspace{0.03cm}\par{\raggedright\tiny\textit{Fuente:} #1\par}%
}
\newcommand{\notafuentes}[1]{\note{[Sources] #1}}
\newcolumntype{Y}{>{\raggedright\arraybackslash}X}

\hypersetup{
  colorlinks=true,
  urlcolor=blue,
  linkcolor=.,
  citecolor=.,
  filecolor=.
}

\newif\ifsectiontoc
\sectiontocfalse
\AtBeginSection[]{%
  \ifsectiontoc
    \begin{frame}{Contenido}
      \tableofcontents[currentsection,hideothersubsections]
    \end{frame}
  \fi
}

\title[Macroeconomía II]{Sesión 5 \\ Modelo intertemporal del sector real con inversión}
\subtitle{Macroeconomía II}
\author[LECO]{Arturo López}
\institute[CIDE]{Centro de Investigación y Docencia Económicas (CIDE)}
\date[Otoño 2026]{Licenciatura en Economía, Otoño 2026}

\begin{document}

\begin{frame}
  \titlepage
\end{frame}

\section{Punto de partida}


\begin{frame}{Objetivos}
  \begin{itemize}
    \setlength{\itemsep}{0.55em}
    \item Integrar las decisiones de consumo, ahorro, trabajo e inversión en un modelo dinámico de equilibrio general.
    \item Construir un modelo que sirva como base para analizar cómo choques macroeconómicos afectan a la economía, y para evaluar el papel de la política macroeconómica.
    \item Un modelo intertemporal del sector real de la economía: mostrar mostrar cómo se determinan la producción real agregada, el consumo real, la inversión real, el empleo, el salario real y la tasa de interés real en la macroeconomía.
    \item Para predecir variables nominales, necesitamos añadir dinero al modelo intertemporal del sector real (próximas sesiones).
    \item El aspecto dinámico del modelo consiste en que hogares y empresas toman decisiones intertemporales: \textit{trade-offs} entre el presente y el futuro.
  \end{itemize}
\end{frame}


\begin{frame}{Decisión intertemporal de la empresa: Inversión}
  \small
  \begin{itemize}
    \item Inversión: gasto en maquinaria, equipo, plantas productivas, etc.
    \item Los \textbf{bienes de inversión} son producidos hoy para su uso futuro en la producción de bienes y servicios.
    \item La capacidad productiva utilizada para producir bienes de inversión podría utilizarse para producir bienes de consumo presente; por otro lado, los bienes de inversión \textbf{incrementan la capacidad productiva futura: el stock de capital futuro}.
    \item Por lo tanto, las empresas hacen un \textbf{\textit{trade-off} entre beneficios presentes y mayores beneficios futuros}.
    \item \textbf{Determinantes de la inversión}: el nivel presente de stock de capital, mayor productividad total de los factores esperada y la tasa de interés real.
    \item Un determinante adicional: percepción de \textbf{riesgo crediticio}. Si los prestamistas (bancos e instituciones financieras) perciben el crédito como más riesgoso en general, las empresas tendrán mayor dificultad en obtener crédito para financiar sus proyectos de inversión.
  \end{itemize}
\end{frame}




\begin{frame}{Setup del modelo}
  \begin{columns}[T,onlytextwidth]
    \begin{column}{0.49\textwidth}
      \textbf{Agentes}
      \begin{itemize}
        \item Hogar representativo.
        \item Empresa representativa.
        \item Gobierno.
      \end{itemize}
    \end{column}
    \begin{column}{0.49\textwidth}
      \textbf{Mercados}
      \begin{itemize}
        \item Laboral, con precio \(w\).
        \item Bienes y servicios, cuyo precio se normaliza a uno.
        \item Crédito, con rendimiento \(r\).
      \end{itemize}
    \end{column}
  \end{columns}
  \vspace{0.4cm}
  \begin{itemize}
    \item Especificamos el modelo en términos de curvas de oferta y demanda.
    \item Podrémos capturar la conducta esencial en el modelo al examinar la \textbf{interacción entre el hogar representativo, la empresa representativa y el gobierno en dos mercados: el mercado laboral en el periodo presente y el merado de bienes y servicios en el periodo presente}.
  \end{itemize}
  \notafuentes{Williamson (2018), pp. 401--402.}
\end{frame}

\begin{frame}{Secuencia económica del modelo}
  \begin{enumerate}
    \setlength{\itemsep}{0.58em}
    \item El hogar elige consumo y ocio en ambos periodos para maximizar su utilidad.
    \item La empresa contrata trabajo, produce el bien de consumo y elige inversión presente.
    \item La inversión determina el capital disponible en el periodo futuro: $K'=(1-\delta)K+I$.
    \item Gobierno, hogar y empresa demandan el bien presente; gobierno y hogar demandan el bien futuro.
    \item \(w, w'\) y \(r\) se ajustan para hacer compatibles todas las decisiones.
  \end{enumerate}
\end{frame}

\section{Hogar representativo}

\begin{frame}{El hogar representativo}
  El hogar dispone de \(h\) unidades de tiempo en cada periodo y elige
  \[
    (C,\ell,C',\ell',S^p).
  \]
  \begin{itemize}
    \setlength{\itemsep}{0.45em}
    \item Oferta laboral: \(N^s=h-\ell\) y \(N^{s\prime}=h-\ell'\).
    \item Ingreso laboral: \(w(h-\ell)\) y \(w'(h-\ell')\).
    \item Recibe beneficios \((\pi,\pi')\) porque es propietario de la empresa.
    \item Paga impuestos de suma fija \((T,T')\).
    \item Toma \((w,w',r)\) como dados.
  \end{itemize}
  \notafuentes{Williamson (2018), pp. 401--402.}
\end{frame}

\begin{frame}{Restricciones presupuestarias por periodo}
  En el presente,
  \[
    C+S^p=w(h-\ell)+\pi-T.
  \]
  En el futuro, el hogar consume todos sus recursos:
  \[
    C'=w'(h-\ell')+\pi'-T'+(1+r)S^p.
  \]
  \begin{itemize}
    \item \(S^p>0\): el hogar es ahorrador.
    \item \(S^p<0\): el hogar es deudor.
  \end{itemize}
  \notafuentes{Williamson (2018), ecuaciones 11-1 y 11-2.}
\end{frame}

\begin{frame}{Restricción presupuestaria intertemporal}
  Eliminando \(S^p\), obtenemos
  \[
    \boxed{
    C+\frac{C'}{1+r}
    =w(h-\ell)+\pi-T
    +\frac{w'(h-\ell')+\pi'-T'}{1+r}
    }.
  \]
  \begin{itemize}
    \setlength{\itemsep}{0.48em}
    \item El lado izquierdo es el valor presente del consumo.
    \item El lado derecho es la \textit{riqueza intertemporal disponible}.
    \item A diferencia del modelo del Cap. 9 (equilibrio parcial), el ingreso disponible por periodo es \textbf{endógeno} (equilibrio general).
  \end{itemize}
  \notafuentes{Williamson (2018), ecuación 11-3.}
\end{frame}

\begin{frame}{Problema del hogar}
  El hogar resuelve
  \[
    \max_{C,C',\ell,\ell'} U(C,\ell,C',\ell') = u(C, \ell) + \beta u(C', \ell')
  \]
  sujeto a
  \[
    C+\frac{C'}{1+r}
    \leq w(h-\ell)+\pi-T
    +\frac{w'(h-\ell')+\pi'-T'}{1+r}.
  \]
  Suponemos que \(U\) es creciente en consumo y ocio, estrictamente cuasicóncava y diferenciable, de modo que la solución interior es única.
  \notafuentes{Williamson (2018), pp. 402--403.}
\end{frame}

\begin{frame}{CPOs}
  \begin{align*}
    &C: \ u_C - \lambda = 0                         \tag{1} \\[4pt]
    &C': \  \beta u_{C'} - \frac{\lambda}{1+r} = 0    \tag{2}  \\[4pt]
    &\ell: \  u_\ell - w\lambda = 0                      \tag{3}  \\[4pt]
    &\ell': \  \beta u_{\ell'} - \frac{w'\lambda}{1+r} = 0 \tag{4} \\[4pt]
    &\lambda: \ w(h-\ell)+\pi-T +\frac{w'(h-\ell')+\pi'-T'}{1+r} - C - \frac{C'}{1+r} = 0 \tag{5}
  \end{align*}
\end{frame}

\begin{frame}{Tres condiciones de optimalidad}
  Sustituyendo (1) en (2):
  $$
  \boxed{u_C = \beta (1+r) u_{C'}} \qquad \text{\textbf{Ecuación de Euler}}
  $$
  Sustituyendo (1) en (3):
  $$
  \boxed{\mathrm{MRS}_{\ell,C}=w} \qquad \text{\textbf{Elección de ocio presente}}
  $$
  Sustituyendo (2) en (4):
  $$
  \boxed{\mathrm{MRS}_{\ell',C'}=w'} \qquad \text{\textbf{Elección de ocio futura}}
  $$
  \begin{itemize}
    \item El salario real presente (futuro) es el costo de oportunidad del ocio presente (futuro).
  \end{itemize}
\end{frame}

\begin{frame}{Oferta laboral en el tiempo}
  \textbf{¿En qué periodo el hogar elegirá trabajar más?}
  
  \vspace{.2cm}

  Sustituyendo (3) en (4):
  \begin{align*}
    \frac{u_\ell}{u_{\ell'}} = \beta(1+r) \frac{w}{w'}
  \end{align*}
  \begin{itemize}
    \small
    \item El ocio será relativamente más alto en el periodo en que el salario sea relativamente más bajo: en periodos de mayor nivel salarial, habrá mayor oferta laboral.
    \item Además, aumentos en $r$ incrementan la oferta laboral: mayor retorno al ahorro $\implies$ mayores incentivos a obtener ingresos laborales $\implies$ mayor ocio y consumo en el futuro.
    \item Esta respuesta de la oferta laboral, motivada por esta \textbf{sustitución intertemporal de trabajo}, está en el corazón de las fluctuaciones de empleo en este modelo.
  \end{itemize}
\end{frame}


\begin{frame}{Determinantes de la oferta laboral presente}
  Las decisiones óptimas generan una función de oferta laboral
  \[
    N^s=N^s(w,r; \omega),
  \]
  donde \(\omega\) denomina la riqueza intertemporal.
  \begin{itemize}
    \setlength{\itemsep}{0.5em}
    \item \(N^s_w>0\): un salario real mayor encarece el costo de oportunidad del ocio (suponemos que el efecto sustitución domina al efecto ingreso).
    \item \(N^s_r>0\): una tasa mayor encarece el ocio presente respecto del ocio futuro (suponemos que el efecto sustitución domina al efecto ingreso).
    \item \(N^s_\omega<0\): aumentar la riqueza intertemporal genera un efecto ingreso.
  \end{itemize}
\end{frame}

\begin{frame}{Respuesta de la oferta laboral presente: $w$}
    \begin{columns}[T,onlytextwidth]
    \begin{column}{0.41\textwidth}
      \small
      \begin{itemize}
        \setlength{\itemsep}{0.48em}
        \item Si \(w\uparrow\), el costo de oportunidad del ocio aumenta.
        \item El hogar sustituye ocio por consumo, suponiendo que domina el efecto sustitución.
        \item Por lo tanto, la curva de oferta laboral tiene pendiente positiva.
      \end{itemize}
    \end{column}
    \begin{column}{0.55\textwidth}
      \centering
      \figura[5.15cm]{fig_11_1_w.png}
      \fuentefigura{Williamson (2018), figura 11.1.}
    \end{column}
  \end{columns}
\end{frame}


\begin{frame}[t]{Respuesta de la oferta laboral presente: $r$}
  \begin{columns}[T,onlytextwidth]
    \begin{column}{0.41\textwidth}
      \small
      \begin{itemize}
        \setlength{\itemsep}{0.48em}
        \item Si \(r\uparrow\), trabajar y ahorrar hoy permite adquirir más consumo y ocio mañana.
        \item El hogar sustituye ocio presente por ocio futuro.
        \item Para todo \(w\), \(N^s\) aumenta.
      \end{itemize}
    \end{column}
    \begin{column}{0.55\textwidth}
      \centering
      \figura[5.15cm]{fig_11_2_oferta_labor_tasa.png}
      \fuentefigura{Williamson (2018), figura 11.2.}
    \end{column}
  \end{columns}
\end{frame}

\begin{frame}{Respuesta de la oferta laboral presente: $\omega$}
    \begin{columns}[T,onlytextwidth]
    \begin{column}{0.41\textwidth}
      \small
      \begin{itemize}
        \setlength{\itemsep}{0.48em}
        \item Si la riqueza intertemporal aumenta \(\omega\uparrow\), se genera un efecto ingreso.
        \item El hogar demanda mayor consumo y ocio (bienes normales).
        \item Para todo \(w\), \(N^s\) cae.
      \end{itemize}
    \end{column}
    \begin{column}{0.55\textwidth}
      \centering
      \figura[5.15cm]{fig_11_3_oferta_laboral_we.png}
      \fuentefigura{Williamson (2018), figura 11.3.}
    \end{column}
  \end{columns}
\end{frame}



\begin{frame}[t]{Demanda de consumo presente}
  \begin{columns}[T,onlytextwidth]
    \begin{column}{0.43\textwidth}
      \small
      \begin{itemize}
        \item La regla de decisión puede escribirse como
        \[
          C^d=C^d(Y,r;\omega).
        \]
        donde $Y$ es el ingreso corriente.
        \item La propensión marginal a consumir satisface
        \[
          0<C_Y^d<1.
        \]
        \item Parte de un aumento del ingreso presente se ahorra (por suavización del consumo).
      \end{itemize}
    \end{column}
    \begin{column}{0.53\textwidth}
      \centering
      \figura[4.7cm]{fig_11_4_consumo_ingreso.png}
      \fuentefigura{Williamson (2018), figura 11.4.}
    \end{column}
  \end{columns}
  \notafuentes{Williamson (2018), figura 11.4 y pp. 406--407.}
\end{frame}

\begin{frame}{Desplazamientos de la demanda de consumo presente}
  Manteniendo constante el ingreso presente \(Y\):
  \begin{itemize}
    \setlength{\itemsep}{0.62em}
    \item Si \(r\uparrow \Rightarrow C^d\downarrow\), suponiendo que domina el efecto sustitución intertemporal.
    \item Si \(Y'\uparrow \Rightarrow C^d\uparrow\), por suavización del consumo.
    \item Si \(\omega \uparrow \Rightarrow C^d\uparrow\), por suavización del consumo.
    \item Por HIP, el desplazamiento de la curva de demanda de consumo debe ser mayor cuando aumenta $\omega$ respecto a cuando aumentan $Y$ o $Y'$.
  \end{itemize}
\end{frame}

\section{Empresa representativa}


\begin{frame}{Tecnología en los dos periodos}
  La empresa produce utilizando capital y trabajo como insumos productivos:
  \[
    Y=zF(K,N),
    \qquad
    Y'=z'F(K',N').
  \]
  \begin{itemize}
    \setlength{\itemsep}{0.55em}
    \item \(z\) y \(z'\) son la productividad total de los factores presente y futura.
    \item \(K\) está predeterminado al inicio del periodo presente.
    \item \(K'\) es elegido indirectamente mediante la inversión actual.
    \item \(F\) tiene productos marginales positivos y decrecientes.
  \end{itemize}
\end{frame}

\begin{frame}{Acumulación de capital}
  La ley de movimiento del capital es
  \[
    \boxed{K'=(1-\delta)K+I},
    \qquad 0\leq\delta\leq1.
  \]
  \begin{itemize}
    \setlength{\itemsep}{0.58em}
    \item Una unidad de producto presente se transforma uno a uno en una unidad de capital futuro.
    \item \((1-\delta)K\) es el capital presente que sobrevive hasta el futuro.
    \item Al final del segundo periodo, el capital remanente se vende como bien de consumo.
  \end{itemize}
  \notafuentes{Williamson (2018), ecuación 11-9.}
\end{frame}

\begin{frame}{Beneficios presentes y futuros}
  Los beneficios son
  \[
    \pi=Y-wN-I,
  \]
  \[
    \pi'=Y'-w'N'+(1-\delta)K'.
  \]
  La empresa maximiza el valor presente de sus beneficios:
  \[
    \boxed{V=\pi+\frac{\pi'}{1+r}}.
  \]
  \begin{itemize}
    \item Invertir reduce los dividendos presentes en una unidad por unidad invertida.
    \item El beneficio aparece mañana mediante mayor producción y mayor valor de liquidación.
  \end{itemize}
  \notafuentes{Williamson (2018), ecuaciones 11-10--11-12.}
\end{frame}

\begin{frame}{Problema de la empresa}
  Tomando \((w,w',r)\) como dados, la empresa resuelve
  \[
    \max_{N,N',I}\;
    V = zF(K,N)-wN-I
    +\frac{z'F((1-\delta)K+I,N')-w'N'+(1-\delta)K'}{1+r}.
  \]
  CPO:
  \begin{align*}
    &N: \ zF_N - w = 0 \tag{1} \\[4pt]
    &N': \ \frac{z'F_{N'} - w'}{1+r} = 0 \tag{2} \\[4pt]
    &I': \ -1 + \frac{z'F_{K'} +(1-\delta)}{1+r} = 0 \tag{3}
  \end{align*}
\end{frame}



\begin{frame}{Condiciones de optimalidad}
  De (1) y (2), las relgas de decisión de demanda laboral son:
  \[
    \boxed{\mathrm{MPN}=w},
    \qquad
    \boxed{\mathrm{MPN}'=w'}.
  \]
\end{frame}


\begin{frame}[t]{Demanda de trabajo presente}
  \begin{columns}[T,onlytextwidth]
    \begin{column}{0.40\textwidth}
      \small
      \begin{itemize}
        \item La empresa contrata trabajo hasta que
        \[
          zF_N(K,N)=w.
        \]
        \item Como \(F_{NN}<0\), la demanda de trabajo tiene pendiente negativa.
        \item \(z\uparrow\) o \(K\uparrow\) desplazan \(N^d\) a la derecha (el $\mathrm{MPN}$ aumenta para todo $N$).
      \end{itemize}
    \end{column}
    \begin{column}{0.56\textwidth}
      \centering
      \figura[5.1cm]{fig_11_7_demanda_labor.png}
      \fuentefigura{Williamson (2018), figura 11.7.}
    \end{column}
  \end{columns}
  \notafuentes{Williamson (2018), figura 11.7 y pp. 410--411.}
\end{frame}


\begin{frame}{Regla óptima de inversión}
  La condición (3) indica que, en el óptimo, el costo marginal en valor presente de la inversión, que es igual a 1, es igual al beneficio marginal en valor presente de la inversión:
  \[
  \boxed{
    \frac{\mathrm{MPK'} + 1-\delta}{1+r} = 1
  }
  \]
  Reordenando la condición de optimalidad, obtenemos la \textbf{regla óptima de inversión}
  \[
  \boxed{
    \mathrm{MPK'} -\delta = r
  }
  \] 
  La empresa invierte hasta el punto en que el \textbf{producto marginal neto del capital}, $\mathrm{MPK'} - \delta$, es igual a la tasa de interés real.
\end{frame}


\begin{frame}{Interpretación de la regla óptima de inversión}
    \begin{itemize}
    \item \(\mathrm{MPK}'-\delta\) es el rendimiento marginal neto del capital.
    \item \(r\) es el costo de oportunidad (el rendimiento disponible en bonos) y el costo efectivo de los proyectos de inversión.
    \item Si \(\mathrm{MPK}'-\delta>r\), la empresa incrementa \(I\).
    \item Si \(\mathrm{MPK}'-\delta<r\), la empresa reduce \(I\).
  \end{itemize}
\end{frame}

\begin{frame}[t]{Demanda de inversión}
  \begin{columns}[T,onlytextwidth]
    \begin{column}{0.40\textwidth}
      \small
      \begin{itemize}
        \item Como \(F_{KK}<0\), un mayor \(K'\) reduce \(\mathrm{MPK}'\).
        \item Dado \(K\), más \(I\) implica más \(K'\).
        \item Por tanto,
        \[
          I^d=I^d(r;K,z',\delta),
          \qquad I_r^d<0.
        \]
      \end{itemize}
    \end{column}
    \begin{column}{0.56\textwidth}
      \centering
      \figura[5.15cm]{fig_11_9_inversion_optima.png}
      \fuentefigura{Williamson (2018), figura 11.9.}
    \end{column}
  \end{columns}
  \notafuentes{Williamson (2018), figura 11.9 y pp. 413--415.}
\end{frame}

\begin{frame}[t]{Desplazamientos de la demanda de inversión}
  \begin{columns}[T,onlytextwidth]
    \begin{column}{0.42\textwidth}
      \small
      \(I^d(r;K,z',\delta)\) se desplaza a la derecha si:
      \begin{itemize}
        \item \(z'\uparrow\): aumenta la productividad futura del capital;
        \item \(K\downarrow\): aumenta el producto marginal futuro del capital;
        \item \(\delta\downarrow\): una fracción mayor del capital sobrevive.
      \end{itemize}
    \end{column}
    \begin{column}{0.54\textwidth}
      \centering
      \figura[4.65cm]{fig_11_10_inversion_desplazamiento.png}
      \fuentefigura{Williamson (2018), figura 11.10.}
    \end{column}
  \end{columns}
  \notafuentes{Williamson (2018), figura 11.10 y pp. 414--416.}
\end{frame}

\begin{frame}{¿Por qué la inversión es tan volátil?}
  \begin{itemize}
    \setlength{\itemsep}{0.63em}
    \item El consumo responde a incentivos de suavización intertemporal.
    \item La inversión responde al rendimiento marginal esperado de proyectos productivos.
    \item Cambios relativamente pequeños en \(r\), \(z'\), \(K\) o el riesgo de crédito pueden generar cambios grandes en \(I\).
    \item La inversión absorbe una parte importante del ajuste ante choques agregados.
  \end{itemize}
  \begin{block}{Predicción}
    La inversión debe ser más volátil que el consumo y, típicamente, más volátil que el producto.
  \end{block}
  \notafuentes{Williamson (2018), pp. 415--416.}
\end{frame}

\begin{frame}{Información asimétrica y prima de incumplimiento}
  Si los bancos no distinguen empresas buenas de malas, las empresas deudoras enfrentan
  \[
    r^L=r+x,
  \]
  donde \(x>0\) es la prima de incumplimiento. Su regla de inversión es
  \[
    \boxed{\mathrm{MPK}'-\delta=r+x}
    \qquad\Longleftrightarrow\qquad
    \mathrm{MPK}'-\delta-x=r.
  \]
  Una mayor incertidumbre crediticia eleva \(x\) y reduce la inversión para cualquier tasa segura \(r\).
  \notafuentes{Williamson (2018), ecuación 11-17 y pp. 417--419.}
\end{frame}

\begin{frame}[t]{Una prima mayor reduce la inversión}
  \begin{columns}[T,onlytextwidth]
    \begin{column}{0.40\textwidth}
      \small
      \begin{itemize}
        \item \(x_2>x_1\) aumenta el costo de financiamiento externo.
        \item La curva relevante se desplaza hacia abajo.
        \item Incluso empresas solventes invierten menos si el mercado no puede identificarlas.
      \end{itemize}
    \end{column}
    \begin{column}{0.56\textwidth}
      \centering
      \figura[5.15cm]{fig_11_11_prima_incumplimiento.png}
      \fuentefigura{Williamson (2018), figura 11.11.}
    \end{column}
  \end{columns}
  \notafuentes{Williamson (2018), figura 11.11 y pp. 418--419.}
\end{frame}

\begin{frame}[t]{Inversión y diferenciales de tasas}
  \begin{columns}[T,onlytextwidth]
    \begin{column}{0.37\textwidth}
      \small
      \begin{itemize}
        \item El diferencial BAA--AAA aproxima la prima de riesgo corporativo.
        \item La inversión y el diferencial exhiben una correlación negativa.
        \item En recesiones, el diferencial tiende a aumentar cuando la inversión cae.
      \end{itemize}
    \end{column}
    \begin{column}{0.59\textwidth}
      \centering
      \figura[5.15cm]{fig_11_12_inversion_spread.png}
      \fuentefigura{Williamson (2018), figura 11.12.}
    \end{column}
  \end{columns}
  \notafuentes{Williamson (2018), figuras 11.12--11.13 y pp. 419--421.}
\end{frame}

\section{Equilibrio}

\begin{frame}{Gobierno}
  El gobierno fija exógenamente \((G,G')\), cobra impuestos de suma fija \((T,T')\) y puede emitir deuda. Su restricción intertemporal es
  \[
    \boxed{G+\frac{G'}{1+r}=T+\frac{T'}{1+r}}.
  \]
  \begin{itemize}
    \setlength{\itemsep}{0.55em}
    \item El valor presente del gasto debe financiarse con el valor presente de impuestos.
    \item Un aumento de \(G\) o \(G'\) reduce la riqueza intertemporal privada.
    \item El momento de los impuestos no altera la riqueza bajo equivalencia ricardiana.
  \end{itemize}
  \notafuentes{Williamson (2018), ecuación 11-18.}
\end{frame}

\begin{frame}{Definición de equilibrio competitivo}
  Dadas las variables exógenas \((K,z,z',G,G')\), un equilibrio competitivo consiste en:
  \begin{enumerate}
    \setlength{\itemsep}{0.48em}
    \item el hogar maximiza su utilidad tomando precios, impuestos y beneficios como dados;
    \item la empresa maximiza el valor presente de sus beneficios tomando precios como dados;
    \item el gobierno satisface su restricción presupuestaria intertemporal;
    \item los mercados de trabajo y de bienes se vacían en ambos periodos.
  \end{enumerate}
  \notafuentes{Williamson (2018), pp. 422--433.}
\end{frame}

\begin{frame}{Condiciones de vaciado de mercado}
  En el periodo presente,
  \[
    \boxed{N=h-\ell},
    \qquad
    \boxed{Y=C+I+G}.
  \]
  En el periodo futuro,
  \[
    \boxed{N'=h-\ell'},
    \qquad
    \boxed{Y'+(1-\delta)K'=C'+G'}.
  \]
  \begin{block}{Ley de Walras}
    Una vez satisfechas las restricciones presupuestarias de hogar, empresa y gobierno, el vaciado de los mercados de bienes y trabajo implica el vaciado del mercado de crédito.
  \end{block}
\end{frame}

\begin{frame}[t]{Mercado laboral para una tasa real dada}
  \begin{columns}[T,onlytextwidth]
    \begin{column}{0.39\textwidth}
      \small
      Para un valor dado de \(r\):
      \begin{enumerate}
        \item \(N^s(w,r)\) y \(N^d(w)\) determinan \((w,N)\).
        \item La función de producción determina
        \[
          Y=zF(K,N).
        \]
        \item Repitiendo para cada \(r\), obtenemos \(Y^s(r)\).
      \end{enumerate}
    \end{column}
    \begin{column}{0.57\textwidth}
      \centering
      \figura[5.2cm]{fig_11_14_labor_produccion.png}
      \fuentefigura{Williamson (2018), figura 11.14.}
    \end{column}
  \end{columns}
  \notafuentes{Williamson (2018), figura 11.14 y pp. 422--424.}
\end{frame}

\begin{frame}{Construcción de la oferta de producto}
  El equilibrio del mercado laboral, dado \(r\), satisface
  \[
    N^s(w,r;\Omega)=N^d(w;K,z).
  \]
  Esta ecuación determina \(w(r)\) y \(N(r)\). Por tanto,
  \[
    \boxed{Y^s(r)=zF(K,N(r))}.
  \]
  Como \(N_r>0\) bajo el supuesto de sustitución intertemporal del ocio,
  \[
    \frac{dY^s}{dr}=zF_N(K,N)\frac{dN}{dr}>0.
  \]
  \notafuentes{Williamson (2018), pp. 423--425.}
\end{frame}

\begin{frame}[t]{La curva de oferta de producto}
  \begin{columns}[T,onlytextwidth]
    \begin{column}{0.39\textwidth}
      \small
      \begin{itemize}
        \item \(r\uparrow\) desplaza la oferta laboral a la derecha.
        \item \(w\downarrow\) y \(N\uparrow\).
        \item Con \(K\) y \(z\) dados, \(Y\uparrow\).
        \item Por ello, \(Y^s\) tiene pendiente positiva en el plano \((Y,r)\).
      \end{itemize}
    \end{column}
    \begin{column}{0.57\textwidth}
      \centering
      \figura[5.25cm]{fig_11_15_oferta_producto.png}
      \fuentefigura{Williamson (2018), figura 11.15.}
    \end{column}
  \end{columns}
  \notafuentes{Williamson (2018), figura 11.15 y pp. 424--425.}
\end{frame}

\begin{frame}{Desplazamientos de la oferta de producto}
  \(Y^s(r)\) se desplaza a la derecha cuando:
  \begin{itemize}
    \setlength{\itemsep}{0.58em}
    \item \(z\uparrow\): aumenta la productividad y la demanda de trabajo;
    \item \(K\uparrow\): aumenta la productividad marginal del trabajo;
    \item \(G\uparrow\) o \(G'\uparrow\): el valor presente de impuestos aumenta, la riqueza privada cae y la oferta laboral sube.
  \end{itemize}
  \begin{alertblock}{Movimiento versus desplazamiento}
    Un cambio en \(r\) genera un movimiento sobre \(Y^s\); cambios en \(z\), \(K\) o la riqueza desplazan la curva.
  \end{alertblock}
  \notafuentes{Williamson (2018), figuras 11.16--11.17 y pp. 425--428.}
\end{frame}

\begin{frame}[t]{Productividad y oferta de producto}
  \begin{columns}[T,onlytextwidth]
    \begin{column}{0.38\textwidth}
      \small
      \begin{itemize}
        \item \(z\uparrow\) desplaza \(N^d\) a la derecha.
        \item La función de producción se desplaza hacia arriba.
        \item Para cada \(r\), aumentan empleo y producto.
        \item \(Y^s\) se desplaza a la derecha.
      \end{itemize}
    \end{column}
    \begin{column}{0.58\textwidth}
      \centering
      \figura[5.25cm]{fig_11_17_ptf_oferta_producto.png}
      \fuentefigura{Williamson (2018), figura 11.17.}
    \end{column}
  \end{columns}
  \notafuentes{Williamson (2018), figura 11.17 y pp. 426--428.}
\end{frame}

\begin{frame}{Demanda de producto: un punto fijo}
  La demanda total de bienes presentes es
  \[
    Y^d=C^d(Y^d,r;\Omega)+I^d(r;K,z',\delta)+G.
  \]
  Para un \(r\) dado, el nivel demandado \(Y^d(r)\) es el punto fijo que satisface
  \[
    \boxed{Y=C^d(Y,r;\Omega)+I^d(r;K,z',\delta)+G}.
  \]
  La condición \(0<C_Y^d<1\) garantiza que el gasto inducido aumenta menos que proporcionalmente con el ingreso.
  \notafuentes{Williamson (2018), ecuación 11-19 y pp. 426--430.}
\end{frame}

\begin{frame}{Pendiente de la demanda de producto}
  Diferenciando implícitamente respecto de \(r\),
  \[
    \frac{dY^d}{dr}
    =\frac{C_r^d+I_r^d}{1-C_Y^d}.
  \]
  Bajo los supuestos
  \[
    C_r^d<0,\qquad I_r^d<0,\qquad 0<C_Y^d<1,
  \]
  se obtiene
  \[
    \boxed{\frac{dY^d}{dr}<0}.
  \]
  Una tasa real mayor reduce consumo e inversión y, mediante el multiplicador, disminuye el producto demandado.
\end{frame}

\begin{frame}[t]{Construcción de la demanda de producto}
  \begin{columns}[T,onlytextwidth]
    \begin{column}{0.38\textwidth}
      \small
      \begin{itemize}
        \item Para cada \(r\), la intersección con la línea de \(45^\circ\) determina \(Y^d(r)\).
        \item \(r_2>r_1\) reduce \(C^d+I^d+G\).
        \item La curva \(Y^d\) tiene pendiente negativa.
      \end{itemize}
    \end{column}
    \begin{column}{0.58\textwidth}
      \centering
      \figura[5.25cm]{fig_11_19_demanda_producto.png}
      \fuentefigura{Williamson (2018), figura 11.19.}
    \end{column}
  \end{columns}
  \notafuentes{Williamson (2018), figura 11.19 y pp. 429--431.}
\end{frame}

\begin{frame}{Desplazamientos de la demanda de producto}
  \(Y^d(r)\) se desplaza a la derecha si:
  \begin{itemize}
    \setlength{\itemsep}{0.54em}
    \item \(G\uparrow\): aumenta directamente la demanda de bienes;
    \item el valor presente de impuestos disminuye: aumenta \(C^d\);
    \item \(Y'\uparrow\): aumenta la riqueza y el consumo presente;
    \item \(z'\uparrow\): aumenta el rendimiento esperado y \(I^d\);
    \item \(K\downarrow\): aumenta \(\mathrm{MPK}'\) y \(I^d\).
  \end{itemize}
  \notafuentes{Williamson (2018), figura 11.20 y pp. 431--432.}
\end{frame}

\begin{frame}{Equilibrio en el mercado de bienes}
  El equilibrio real intertemporal satisface
  \[
    \boxed{Y^s(r^*)=Y^d(r^*)}.
  \]
  La intersección determina simultáneamente
  \[
    (Y^*,r^*).
  \]
  Dado \(r^*\), el mercado laboral determina
  \[
    (N^*,w^*).
  \]
  Finalmente, las reglas de decisión de los agentes determinan \(C^*\), \(I^*\), \(\ell^*\) y \(K^{\prime *}\).
\end{frame}

\begin{frame}[t]{El modelo intertemporal real completo}
  \begin{columns}[T,onlytextwidth]
    \begin{column}{0.36\textwidth}
      \small
      \begin{itemize}
        \item Panel (b): \(Y^d=Y^s\) determina \(r^*\) y \(Y^*\).
        \item Panel (a): dado \(r^*\), \(N^d=N^s\) determina \(w^*\) y \(N^*\).
        \item Los dos mercados se resuelven simultáneamente.
      \end{itemize}
    \end{column}
    \begin{column}{0.60\textwidth}
      \centering
      \figura[4.95cm]{fig_11_21_modelo_completo.png}
      \fuentefigura{Williamson (2018), figura 11.21.}
    \end{column}
  \end{columns}
  \notafuentes{Williamson (2018), figura 11.21 y pp. 432--434.}
\end{frame}

\begin{frame}{Orden práctico para resolver el modelo}
  \begin{enumerate}
    \setlength{\itemsep}{0.62em}
    \item Derivar \(N^s(w,r)\), \(C^d(Y,r)\), \(N^d(w)\) e \(I^d(r)\).
    \item Imponer \(N^s=N^d\) para construir \(Y^s(r)\).
    \item Imponer \(Y=C^d+I^d+G\) para construir \(Y^d(r)\).
    \item Resolver \(Y^s(r^*)=Y^d(r^*)\).
    \item Recuperar \(Y^*,w^*,N^*,C^*,I^*\) y \(K^{\prime *}\).
  \end{enumerate}
  \begin{block}{Experimentación}
    Un choque modifica una o más reglas de decisión; esas modificaciones desplazan curvas y producen un nuevo vector de precios y asignaciones.
  \end{block}
\end{frame}

\section{Experimentos}

\begin{frame}{Cómo leer un choque en equilibrio general}
  \begin{enumerate}
    \setlength{\itemsep}{0.55em}
    \item Identificar la variable exógena que cambia y el mercado donde aparece directamente.
    \item Determinar qué reglas de decisión se modifican.
    \item Traducir esos cambios en desplazamientos de \(Y^s(r)\), \(Y^d(r)\) o de las curvas del mercado laboral.
    \item Encontrar el nuevo \((Y^*,r^*)\) y, después, el nuevo \((N^*,w^*)\).
    \item Seguir la propagación del choque hacia consumo, inversión y capital futuro.
  \end{enumerate}
  \begin{alertblock}{Disciplina del análisis}
    Un movimiento sobre una curva proviene del ajuste de un precio endógeno; un desplazamiento proviene de un cambio exógeno o paramétrico.
  \end{alertblock}
\end{frame}

\begin{frame}{Aumento temporal del gasto público}
  Suponga que \(G\) aumenta en el periodo presente, mientras \(G'\) permanece constante.
  \begin{itemize}
    \setlength{\itemsep}{0.48em}
    \item El choque se origina en el mercado de bienes: la demanda pública aumenta directamente.
    \item El mayor valor presente de impuestos reduce la riqueza del hogar y desplaza la demanda de consumo hacia abajo.
    \item La menor riqueza incrementa la oferta laboral y desplaza la oferta de producto hacia la derecha.
    \item La inversión no se desplaza directamente: se mueve sobre \(I^d(r)\) cuando cambia \(r\).
  \end{itemize}
  \[
    \Delta Y^d
    =
    \Delta G+\Delta C^d,
    \qquad
    \Delta C^d<0.
  \]
\end{frame}

\begin{frame}{Desplazamiento inicial de la demanda de producto}
  Sea \(m\in(0,1)\) la propensión marginal a consumir respecto de la riqueza intertemporal y sea \(\Delta\) el desplazamiento horizontal de \(Y^d\), manteniendo \(r\) fijo. Entonces,
  \[
    \Delta
    =
    \Delta G-m\Delta G+m\Delta.
  \]
  Por lo tanto,
  \[
    (1-m)\Delta=(1-m)\Delta G
    \quad\Longrightarrow\quad
    \boxed{\Delta=\Delta G}.
  \]
  \begin{block}{Interpretación}
    A una tasa de interés dada, el aumento del gasto desplaza la demanda de producto por el monto de \(\Delta G\); el ajuste endógeno de la producción genera ingreso adicional y consumo inducido.
  \end{block}
\end{frame}

\begin{frame}[t]{Equilibrio ante un aumento temporal de \(G\)}
  \begin{columns}[T,onlytextwidth]
    \begin{column}{0.38\textwidth}
      \small
      \begin{itemize}
        \item \(Y^d\) se desplaza a la derecha.
        \item \(Y^s\) también se desplaza a la derecha por el efecto riqueza sobre la oferta laboral.
        \item En el resultado mostrado por Williamson:
        \[
          Y\uparrow,\quad r\uparrow,\quad
          N\uparrow,\quad w\downarrow.
        \]
        \item El mayor \(r\) desplaza gasto privado:
        \[
          C\downarrow,\qquad I\downarrow.
        \]
      \end{itemize}
    \end{column}
    \begin{column}{0.58\textwidth}
      \centering
      \figura[5.05cm]{fig_11_22_gasto_publico.png}
      \fuentefigura{Williamson (2018), figura 11.22.}
    \end{column}
  \end{columns}
  \notafuentes{Williamson (2018), figura 11.22 y pp. 434--437.}
\end{frame}

\begin{frame}{Crowding out y multiplicador del gasto}
  Del vaciado del mercado de bienes,
  \[
    \Delta Y=\Delta C+\Delta I+\Delta G.
  \]
  Como el aumento de \(G\) eleva la tasa de interés y reduce riqueza,
  \[
    \Delta C<0,
    \qquad
    \Delta I<0.
  \]
  Por consiguiente,
  \[
    0<\frac{\Delta Y}{\Delta G}<1
  \]
  en la configuración estudiada por Williamson.
  \begin{block}{Mecanismo}
    Parte del gasto público adicional desplaza consumo e inversión privados. El producto aumenta, pero menos que proporcionalmente al impulso fiscal.
  \end{block}
\end{frame}

\begin{frame}[t]{Gasto público, producto y consumo en los datos}
  \begin{columns}[T,onlytextwidth]
    \begin{column}{0.34\textwidth}
      \small
      \begin{itemize}
        \item La evidencia descriptiva no identifica por sí sola un efecto causal.
        \item El gasto puede responder endógenamente al ciclo económico.
        \item El modelo disciplina qué comovimientos buscar después de un choque fiscal exógeno.
      \end{itemize}
    \end{column}
    \begin{column}{0.62\textwidth}
      \centering
      \figura[5.15cm]{fig_11_23_gasto_pib_consumo.png}
      \fuentefigura{Williamson (2018), figura 11.23.}
    \end{column}
  \end{columns}
  \notafuentes{Williamson (2018), figura 11.23 y pp. 437--438.}
\end{frame}

\begin{frame}{Destrucción del acervo de capital}
  Una caída exógena de \(K\) opera por dos canales:
  \begin{enumerate}
    \setlength{\itemsep}{0.52em}
    \item \textbf{Oferta presente:} cae la productividad marginal del trabajo, de modo que \(N^d\) y \(Y^s\) se desplazan a la izquierda.
    \item \textbf{Inversión:} aumenta el producto marginal futuro del capital, por lo que \(I^d\) y \(Y^d\) se desplazan a la derecha.
  \end{enumerate}
  En el nuevo equilibrio,
  \[
    Y\,?,\qquad r\uparrow,\qquad
    C\downarrow,\qquad I\uparrow,\qquad w\downarrow.
  \]
  Los efectos sobre producto y empleo son ambiguos. En la figura, dominan los desplazamientos de oferta y ambas variables disminuyen.
\end{frame}

\begin{frame}[t]{Equilibrio después de una pérdida de capital}
  \begin{columns}[T,onlytextwidth]
    \begin{column}{0.37\textwidth}
      \small
      \begin{itemize}
        \item La figura muestra el caso en que la economía produce menos hoy.
        \item La escasez de capital hace rentable reconstruirlo.
        \item La tasa real aumenta para reasignar bienes desde consumo hacia inversión.
        \item En general, el producto puede aumentar o disminuir, según cuál desplazamiento domine.
      \end{itemize}
    \end{column}
    \begin{column}{0.59\textwidth}
      \centering
      \figura[5.05cm]{fig_11_24_capital.png}
      \fuentefigura{Williamson (2018), figura 11.24.}
    \end{column}
  \end{columns}
  \notafuentes{Williamson (2018), figura 11.24 y pp. 438--440.}
\end{frame}

\begin{frame}{Aumento de la productividad presente}
  Considere un aumento transitorio de \(z\), manteniendo \(z'\) constante.
  \begin{itemize}
    \setlength{\itemsep}{0.52em}
    \item La productividad marginal del trabajo aumenta: \(N^d\) y \(Y^s\) se desplazan a la derecha.
    \item La curva \(Y^d(r)\) no se desplaza directamente.
    \item El nuevo equilibrio se alcanza mediante un movimiento sobre \(Y^d(r)\): la tasa real disminuye.
    \item La menor tasa real eleva consumo e inversión; el mayor ingreso corriente refuerza el aumento del consumo.
  \end{itemize}
  En el caso de referencia,
  \[
    Y\uparrow,\quad r\downarrow,\quad
    C\uparrow,\quad I\uparrow,\quad w\uparrow.
  \]
\end{frame}

\begin{frame}[t]{Equilibrio ante mayor productividad presente}
  \begin{columns}[T,onlytextwidth]
    \begin{column}{0.37\textwidth}
      \small
      \begin{itemize}
        \item El choque se origina en la tecnología de producción presente.
        \item La oferta de producto se desplaza más que la demanda.
        \item El efecto sobre \(N\) combina una mayor demanda laboral con efectos ingreso y sustitución sobre la oferta.
        \item Williamson obtiene \(N\uparrow\) bajo su configuración gráfica.
      \end{itemize}
    \end{column}
    \begin{column}{0.59\textwidth}
      \centering
      \figura[5.05cm]{fig_11_25_ptf_actual.png}
      \fuentefigura{Williamson (2018), figura 11.25.}
    \end{column}
  \end{columns}
  \notafuentes{Williamson (2018), figura 11.25 y pp. 440--442.}
\end{frame}

\begin{frame}[t]{Productividad laboral media y ciclo económico}
  \begin{columns}[T,onlytextwidth]
    \begin{column}{0.34\textwidth}
      \small
      \begin{itemize}
        \item En el modelo, un choque positivo de productividad aumenta producto y productividad laboral.
        \item Esto genera comovimiento procíclico entre ambas variables.
        \item Sin embargo, la productividad observada también refleja composición sectorial, utilización del capital y medición.
      \end{itemize}
    \end{column}
    \begin{column}{0.62\textwidth}
      \centering
      \figura[5.15cm]{fig_11_26_productividad_media.png}
      \fuentefigura{Williamson (2018), figura 11.26.}
    \end{column}
  \end{columns}
  \notafuentes{Williamson (2018), figura 11.26 y pp. 442--443.}
\end{frame}

\begin{frame}{Noticias sobre productividad futura}
  Suponga que aumenta \(z'\), sin cambio en \(z\).
  \begin{itemize}
    \setlength{\itemsep}{0.5em}
    \item El mayor producto marginal futuro del capital desplaza \(I^d(r)\) y \(Y^d(r)\) a la derecha.
    \item La tasa real aumenta para equilibrar la mayor demanda de bienes.
    \item La mayor tasa real desplaza trabajo hacia el presente: la oferta laboral aumenta.
  \end{itemize}
  El ajuste de equilibrio requiere
  \[
    r\uparrow,\qquad Y\uparrow,\qquad N\uparrow,\qquad
    w\downarrow,\qquad I\uparrow.
  \]
  El efecto sobre \(C\) es, en general, ambiguo porque el efecto riqueza positivo compite con el mayor rendimiento del ahorro.
\end{frame}

\begin{frame}[t]{Equilibrio ante mejores perspectivas futuras}
  \begin{columns}[T,onlytextwidth]
    \begin{column}{0.37\textwidth}
      \small
      \begin{itemize}
        \item \(Y^d\) se desplaza a la derecha por inversión y consumo.
        \item \(Y^s(r)\) no se desplaza directamente; la economía se mueve sobre esta curva.
        \item La tasa real aumenta para contener demanda presente y movilizar recursos hacia inversión.
        \item El mayor \(r\) aumenta la oferta laboral y reduce el salario real.
      \end{itemize}
    \end{column}
    \begin{column}{0.59\textwidth}
      \centering
      \figura[5.05cm]{fig_11_27_ptf_futura.png}
      \fuentefigura{Williamson (2018), figura 11.27.}
    \end{column}
  \end{columns}
  \notafuentes{Williamson (2018), figura 11.27 y pp. 443--446.}
\end{frame}

\begin{frame}[t]{Inversión y mercado accionario}
  \begin{columns}[T,onlytextwidth]
    \begin{column}{0.35\textwidth}
      \small
      \begin{itemize}
        \item Los precios accionarios incorporan expectativas sobre beneficios futuros.
        \item Noticias favorables pueden elevar simultáneamente la valuación de las empresas y la inversión.
        \item La correlación no prueba que la bolsa cause la inversión: ambas pueden responder a la misma información.
      \end{itemize}
    \end{column}
    \begin{column}{0.61\textwidth}
      \centering
      \figura[5.10cm]{fig_11_28_inversion_bolsa.png}
      \fuentefigura{Williamson (2018), figura 11.28.}
    \end{column}
  \end{columns}
  \notafuentes{Williamson (2018), figura 11.28 y pp. 446--448.}
\end{frame}

\begin{frame}{Aumento de las fricciones financieras}
  Una mayor prima de incumplimiento \(d\) eleva el costo relevante para financiar inversión:
  \[
    1+r+d.
  \]
  \begin{itemize}
    \setlength{\itemsep}{0.5em}
    \item \(I^d(r;d)\) se desplaza a la izquierda.
    \item La menor inversión reduce \(Y^d(r)\).
    \item Las restricciones de crédito y la mayor tasa efectiva para los deudores inducen más oferta laboral; \(Y^s(r)\) se desplaza a la derecha.
    \item La tasa real de mercado disminuye, pero no compensa completamente la prima.
  \end{itemize}
  El libro obtiene \(Y\downarrow\), \(N\downarrow\), \(r\downarrow\) y \(w\uparrow\); la composición entre \(C\) e \(I\) depende de parámetros.
\end{frame}

\begin{frame}[t]{Equilibrio con mayores fricciones de crédito}
  \begin{columns}[T,onlytextwidth]
    \begin{column}{0.37\textwidth}
      \small
      \begin{itemize}
        \item El choque se origina en el mercado de crédito.
        \item Se encarece la transformación de bienes presentes en capital futuro.
        \item La inversión cae y la menor demanda se transmite al producto y al empleo.
        \item Éste es un mecanismo central para interpretar crisis financieras.
      \end{itemize}
    \end{column}
    \begin{column}{0.59\textwidth}
      \centering
      \figura[5.05cm]{fig_11_29_fricciones_credito.png}
      \fuentefigura{Williamson (2018), figura 11.29.}
    \end{column}
  \end{columns}
  \notafuentes{Williamson (2018), figura 11.29 y pp. 448--451.}
\end{frame}

\begin{frame}{Choque sectorial y reasignación}
  Considere una perturbación que destruye emparejamientos productivos entre trabajadores y empresas.
  \begin{itemize}
    \setlength{\itemsep}{0.5em}
    \item La productividad agregada efectiva disminuye porque parte del trabajo y del capital queda temporalmente mal asignada.
    \item La oferta de producto se desplaza a la izquierda.
    \item El desempleo puede aumentar aun cuando la productividad media observada aumente, si primero desaparecen los empleos menos productivos.
    \item La reasignación toma tiempo, por lo que el choque puede tener efectos persistentes.
  \end{itemize}
  \[
    Y\downarrow,\qquad N\downarrow,\qquad r\uparrow,
  \]
  mientras que el efecto sobre \(w\) depende de la composición y de la intensidad de los desplazamientos.
\end{frame}

\begin{frame}[t]{Equilibrio ante un choque de reasignación}
  \begin{columns}[T,onlytextwidth]
    \begin{column}{0.37\textwidth}
      \small
      \begin{itemize}
        \item La capacidad productiva corriente cae.
        \item El producto y el empleo disminuyen.
        \item La tasa real aumenta para reducir la demanda presente.
        \item El modelo permite una recesión con productividad media creciente.
      \end{itemize}
    \end{column}
    \begin{column}{0.59\textwidth}
      \centering
      \figura[5.05cm]{fig_11_30_choque_sectorial.png}
      \fuentefigura{Williamson (2018), figura 11.30.}
    \end{column}
  \end{columns}
  \notafuentes{Williamson (2018), figura 11.30 y pp. 451--454.}
\end{frame}

\begin{frame}[t]{Dos recesiones, distintos mecanismos}
  \begin{columns}[T,onlytextwidth]
    \begin{column}{0.49\textwidth}
      \centering
      \textbf{Producto real}

      \figura[3.85cm]{fig_11_32_pib_recesiones.png}
    \end{column}
    \begin{column}{0.49\textwidth}
      \centering
      \textbf{Empleo}

      \figura[3.85cm]{fig_11_33_empleo_recesiones.png}
    \end{column}
  \end{columns}
  \begin{itemize}
    \small
    \item Tanto la recesión de 1981--1982 como la de 2008--2009 redujeron producto y empleo, pero la composición y persistencia de los movimientos fueron diferentes.
  \end{itemize}
  \notafuentes{Williamson (2018), figuras 11.32 y 11.33.}
\end{frame}

\begin{frame}[t]{Productividad durante dos recesiones}
  \begin{columns}[T,onlytextwidth]
    \begin{column}{0.35\textwidth}
      \small
      \begin{itemize}
        \item La productividad se comportó de manera distinta en ambas recesiones.
        \item Una sola perturbación tecnológica no explica necesariamente todos los episodios.
        \item El diagnóstico requiere confrontar predicciones conjuntas sobre \(Y\), \(N\), productividad, \(C\), \(I\) y \(r\).
      \end{itemize}
    \end{column}
    \begin{column}{0.61\textwidth}
      \centering
      \figura[5.10cm]{fig_11_34_productividad_recesiones.png}
      \fuentefigura{Williamson (2018), figura 11.34.}
    \end{column}
  \end{columns}
  \notafuentes{Williamson (2018), figura 11.34 y pp. 454--456.}
\end{frame}

\section{Síntesis}

\begin{frame}{Predicciones comparativas del modelo}
  \small
  \begin{itemize}
    \setlength{\itemsep}{0.55em}
    \item \(G\) temporal \(\uparrow\):
    \(Y\uparrow,\ r\uparrow,\ N\uparrow,\ w\downarrow,\ C\downarrow,\ I\downarrow\).
    \item \(K\downarrow\):
    \(Y\,?,\ r\uparrow,\ N\,?,\ w\downarrow,\ C\downarrow,\ I\uparrow\).
    \item \(z\uparrow\):
    \(Y\uparrow,\ r\downarrow,\ N\,?,\ w\uparrow,\ C\uparrow,\ I\uparrow\).
    \item \(z'\uparrow\):
    \(Y\uparrow,\ r\uparrow,\ N\uparrow,\ w\downarrow,\ C\,?,\ I\uparrow\).
    \item Fricción financiera \(\uparrow\):
    \(Y\downarrow,\ r\downarrow,\ N\downarrow,\ w\uparrow,\ C\,?,\ I\,?\).
    \item Reasignación adversa:
    \(Y\downarrow,\ r\uparrow,\ N\downarrow,\ w\,?,\ C\downarrow,\ I\downarrow\).
  \end{itemize}
  \footnotesize Los signos corresponden a la configuración analizada por Williamson; \(?\) indica ambigüedad sin restricciones adicionales.
\end{frame}

\begin{frame}{Resultados centrales}
  \begin{itemize}
    \setlength{\itemsep}{0.58em}
    \item La inversión conecta el equilibrio presente con la capacidad productiva futura.
    \item La tasa de interés real coordina ahorro e inversión y equilibra la demanda con la oferta de producto.
    \item La decisión de invertir iguala el costo marginal presente con el valor descontado del producto marginal futuro del capital.
    \item Choques fiscales, tecnológicos, financieros y de reasignación se distinguen por sus predicciones conjuntas, no por el movimiento de una sola variable.
    \item El modelo ofrece un lenguaje disciplinado para identificar el origen, la propagación y los mecanismos de transmisión de las fluctuaciones macroeconómicas.
  \end{itemize}
\end{frame}

\begin{frame}{Referencias}
  \footnotesize
  \begin{thebibliography}{9}
    \bibitem{Williamson2018}
    Williamson, S. D. (2018).
    \newblock \textit{Macroeconomics}, 6a ed.
    \newblock Pearson, capítulo 11: \textit{A Real Intertemporal Model with Investment}.
  \end{thebibliography}
\end{frame}

\appendix
\section{Apéndice}

\begin{frame}{Apéndice: derivación de la inversión óptima}
  Sustituyendo \(K'=(1-\delta)K+I\), el componente de beneficios que depende de \(I\) es
  \[
    -I+\frac{z'F((1-\delta)K+I,N')+(1-\delta)K'}{1+r}.
  \]
  La condición de primer orden es
  \[
    -1+\frac{z'F_K(K',N')+(1-\delta)}{1+r}=0.
  \]
  Por lo tanto,
  \[
    \boxed{1+r=z'F_K(K',N')+1-\delta}
  \]
  o, equivalentemente,
  \[
    \boxed{r=z'F_K(K',N')-\delta}.
  \]
\end{frame}

\begin{frame}{Apéndice: crédito y ley de Walras}
  El ahorro del hogar financia la inversión neta de la empresa y la posición fiscal del gobierno. Al agregar las restricciones presupuestarias privadas y utilizar el presupuesto público, se obtiene
  \[
    Y-C-G-I=0.
  \]
  Así, una vez que se satisfacen las restricciones presupuestarias y se vacían los mercados de trabajo y crédito, el mercado de bienes también se vacía.
  \begin{block}{Implicación}
    No todas las condiciones de vaciado son independientes. La ley de Walras permite omitir una de ellas al resolver el sistema, pero no elimina la obligación de definirla y verificarla.
  \end{block}
\end{frame}

\begin{frame}{Apéndice: interpretación mediante la \(q\) de Tobin}
  Defina el valor marginal de una unidad adicional de capital instalado como
  \[
    q
    \equiv
    \frac{z'F_K(K',N')+1-\delta}{1+r}.
  \]
  Sin costos de ajuste, la condición óptima de inversión es
  \[
    q=1.
  \]
  \begin{itemize}
    \item Si \(q>1\), el valor de una unidad adicional de capital excede su costo: conviene invertir más.
    \item Si \(q<1\), conviene reducir la inversión.
    \item Con costos convexos de ajuste, la inversión responde gradualmente y \(q\) puede diferir de uno en equilibrio.
  \end{itemize}
\end{frame}

\begin{frame}{Apéndice: alcances y límites}
  \begin{itemize}
    \setlength{\itemsep}{0.52em}
    \item El modelo tiene dos periodos y certidumbre; no describe expectativas estocásticas ni dinámica de largo plazo.
    \item La empresa puede transformar inversión en capital sin costos de ajuste ni rezagos de instalación.
    \item Los mercados son competitivos salvo por la prima financiera introducida de forma reducida.
    \item No hay rigideces nominales, dinero ni política monetaria.
    \item A pesar de estas simplificaciones, el modelo contiene los márgenes esenciales de consumo, trabajo, ahorro, inversión y acumulación de capital.
  \end{itemize}
\end{frame}

\end{document}
